\paragraph{RGB visuomotor policies.}
Modern imitation-learning policies map image observations and proprioception to
action chunks with expressive generative decoders, including diffusion and
flow-matching action experts~\citep{chi2023diffusionpolicy,yan2025maniflow}.
These policies are strong action predictors but treat each image as raw
appearance: camera geometry is implicit, so a fixed encoder can exploit
view-specific cues that need not transfer when camera placement changes.  Our
work keeps the RGB-policy interface but makes camera pose an explicit, learned
latent inside the policy.

\paragraph{Policies with 3D and geometric structure.}
A complementary line builds geometric structure into the policy.  Point-cloud
policies consume depth-lifted 3D inputs with modality-specific encoders: 3D
Diffusion Policy~\citep{ze2024dp3} shows that compact 3D representations improve
sample efficiency, and ManiFlow~\citep{yan2025maniflow} scales a point-cloud
flow-matching policy across many tasks.  Other methods inject geometric priors
through structure rather than explicit 3D inputs: Equivariant Diffusion
Policy~\citep{wang2024equidiff} bakes $SO(2)/SO(3)$ symmetry into the denoiser
for better sample efficiency and generalization, and Hierarchical Equivariant
Policy~\citep{zhao2025hep} composes equivariant sub-policies through a learned
frame-transfer hierarchy.  These approaches demonstrate the value of geometric
structure, but they require depth and calibration (point-cloud methods) or
assume a specific symmetry or kinematic decomposition.  We instead recover
camera geometry from RGB inside the control objective, so the policy can still
train on episodes without depth or extrinsic labels and reuse standard image
backbones.

\paragraph{Camera conditioning for RGB policies.}
Recent camera-aware RGB policies expose geometry to the encoder through
calibrated camera encodings.  Per-pixel Plucker ray maps concatenate ray origin
and direction to RGB tokens~\citep{sitzmann2021lfn}, and attention-level
encodings such as GTA and PRoPE inject relative camera pose inside
self-attention~\citep{miyato2024gta,li2025prope}; calibrated multi-view policies
use these signals for viewpoint
generalization~\citep{anonymous2026cameraaware}.  All of these supply geometry
as fixed metadata and assume known extrinsics at both training and test time.
Our policies use the same Plucker ray encoding but build it from \emph{predicted}
pose, removing the inference-time calibration requirement.

\paragraph{View synthesis and camera-pose estimation.}
Ray- and pose-conditioned transformers underpin recent view-synthesis models:
set-latent scene representations~\citep{sajjadi2022srt}, large geometry-free
view-synthesis transformers~\citep{jin2024lvsm}, and self-supervised,
unposed-and-uncalibrated variants that predict cameras as a by-product of
rendering~\citep{jiang2025rayzer}.  Pose can also be regressed directly from
images, e.g.\ as ray bundles~\citep{zhang2024camerasasrays}.  We borrow the
ray-map conditioning and set-latent scene tokens from this literature, but use
target-view reconstruction only as a training-time regularizer that forces the
scene representation to stay geometrically faithful for a control policy, rather
than as the end task.
